% interaction/security.tex — completeness, soundness, knowledge soundness, claim trees

\section{Security Definitions}\label{sec:interaction-security}

Security definitions are conceptually the same as in the flat-indexed model but
now sit on the $\mathsf{Spec} + \mathsf{RoleDecoration}$ foundation.  At the
oracle level, the key difference is that verifier-side acceptance is phrased in
terms of the \emph{existence} of concrete oracle data compatible with the
verifier's query-level simulation, rather than by assuming a built-in
reification function.

Thus the security layer changes much less than the interaction layer itself:
the semantics are richer, but the core completeness/soundness/knowledge
soundness pattern remains the same.

\subsection{Protocol participants}

\begin{definition}[Prover, Verifier, Reduction]
  \label{int:reduction}
  A \emph{prover} takes $(\mathit{stmt}, \mathit{wit})$ and produces, via
  monadic setup, a $\Strategy.\mathsf{withRoles}$ whose output is
  $\mathsf{HonestProverOutput}\;(\mathit{StmtOut}\;s\;\mathit{tr})\;
  (\mathit{WitOut}\;s\;\mathit{tr})$.
  A \emph{verifier} is a statement-indexed $\mathsf{Counterpart}$ with
  $\mathit{StmtOut}\;s\;\mathit{tr}$ at $\mathsf{done}$.
  A \emph{reduction} pairs a prover with a verifier.
  $\mathit{WitnessIn}$ is intentionally not statement-dependent; compatibility
  is expressed in the security relations.
  \lean{Interaction.Reduction}
  \uses{int:strategy-with-roles, int:counterpart}
\end{definition}

$\Reduction.\mathsf{execute}$ runs the prover's strategy against the verifier
via $\mathsf{runWithRoles}$, returning the transcript and both outputs.
$\Reduction.\mathsf{Continuation}$ supports transcript-indexed second-stage
composition.

\subsection{Completeness, soundness, knowledge soundness}

All definitions use a generic monad~$m$ with $[\mathsf{HasEvalSPMF}\;m]$ for
probability semantics.

\begin{definition}[Completeness]
  \label{int:completeness}
  A reduction satisfies \emph{completeness with error $\varepsilon$} if for all
  valid inputs, honest execution produces valid output with probability at
  least $1 - \varepsilon$.
  \lean{Interaction.Reduction.completeness}
  \uses{int:reduction}
\end{definition}

\begin{definition}[Soundness]
  \label{int:soundness}
  A verifier satisfies \emph{soundness with error $\varepsilon$} if for all
  inputs outside the input language and all (possibly malicious) provers, the
  probability that the verifier's output falls in the output language is at
  most~$\varepsilon$.  Thus soundness is verifier-side: the honest prover does
  not appear in the definition.
  \lean{Interaction.soundness}
  \uses{int:reduction}
\end{definition}

\begin{definition}[Knowledge Soundness]
  \label{int:knowledge-soundness}
  A verifier satisfies \emph{knowledge soundness with error $\varepsilon$} if
  there exists a straightline extractor such that for all provers, the
  probability that the verifier accepts but the extractor fails to produce a
  valid witness is at most~$\varepsilon$.
  \lean{Interaction.knowledgeSoundness}
  \uses{int:reduction}
\end{definition}

\subsection{Composition theorems}

\begin{theorem}[Completeness composes]
  \label{thm:completeness-comp}
  If reduction~1 has completeness error~$\varepsilon_1$ and reduction~2
  (a continuation) has completeness error~$\varepsilon_2$ at every first-phase
  transcript, then the composed reduction has completeness error at most
  $\varepsilon_1 + \varepsilon_2$.
  \lean{Interaction.Reduction.completeness_comp}
  \uses{int:completeness}
\end{theorem}

Analogous theorems hold for perfect completeness and soundness.

\subsection{Round-by-round analysis}

\begin{definition}[Claim Tree]
  \label{int:claim-tree}
  A \emph{claim tree} is a recursive soundness witness on $\mathsf{Spec} +
  \mathsf{RoleDecoration}$.  At each sender (prover-message) node, bad claims
  must stay bad.  At each receiver (verifier-challenge) node, a bad claim may
  flip to good with probability at most $\mathit{error}$.
  \lean{Interaction.ClaimTree}
  \uses{int:role-decoration}
\end{definition}

\begin{theorem}[Terminal probability bound]
  \label{thm:bound-terminal-prob}
  If a $\mathsf{ClaimTree}$ is sound with per-round error bounds, then the
  probability of reaching a good terminal claim from a bad root is bounded by
  the sum of the per-round errors along any path.
  \lean{Interaction.ClaimTree.IsSound.bound_terminalProb}
  \uses{int:claim-tree}
\end{theorem}

\subsection{Oracle security}

At the oracle level, verifier-side acceptance is phrased via \emph{existence}
of concrete output oracle data:

\begin{definition}[Oracle Completeness]
  \label{int:oracle-completeness}
  Honest execution produces concrete output oracle data that realizes the
  verifier's simulation, and the output passes the acceptance predicate.
  \lean{Interaction.OracleDecoration.OracleReduction.completeness}
  \uses{int:oracle-reduction, int:completeness}
\end{definition}

\begin{definition}[Oracle Soundness]
  \label{int:oracle-soundness}
  For any malicious prover, the probability that there exists a concrete output
  oracle family realizing the verifier's simulation \emph{and} the resulting
  output passes acceptance is at most~$\varepsilon$.
  \lean{Interaction.OracleDecoration.OracleReduction.soundness}
  \uses{int:oracle-reduction, int:soundness}
\end{definition}
